\chapter{Future Scope}
\setcounter{equation}{0}
\section{Directions Discussed in the Paper}
The paper identifies opportunities for improving deployment efficiency and adapting to changing URL patterns. These include model compression, automatic retraining, multi-model ensembles, and runtime validation using checks such as landing-page rendering and WHOIS/DNS lookups alongside static model predictions.\cite{bib1}

\section{Possible Extensions}
\begin{itemize}
\item \textbf{Model compression:} Investigate tree pruning and export formats such as ONNX or Treelite, including quantisation where appropriate, to reduce runtime and storage requirements.
\item \textbf{Automated model refresh:} Develop a monitored retraining pipeline using newly verified phishing and legitimate URLs, with safeguards against poisoned labels.
\item \textbf{Ensemble detection:} Evaluate combinations of gradient-boosted trees, neural methods, and other classifiers to determine whether complementary signals improve robust performance.
\item \textbf{Runtime checks:} Combine URL features with carefully selected live DNS, WHOIS, redirect, and page-level signals, while measuring the extra latency and privacy implications.
\item \textbf{Broader browser support:} Test performance and usability across browsers and operating environments.
\item \textbf{Continuous evaluation:} Report precision, recall, F1, false-positive rate, latency, and memory on temporally separated and independent datasets.
\end{itemize}

These are potential research and engineering extensions; they should not be interpreted as improvements already validated by the paper.
